% ============================================================================
% Bang dien giai tu viet tat va tieng Anh
% ============================================================================

\chapter*{BẢNG DIỄN GIẢI TỪ VIẾT TẮT VÀ TIẾNG ANH}
\addcontentsline{toc}{chapter}{BẢNG DIỄN GIẢI TỪ VIẾT TẮT VÀ TIẾNG ANH}

% Sap xep theo thu tu ABC cua chu cai dau tien

\begin{longtable}{@{}p{3.5cm} >{\raggedright\arraybackslash}p{6cm} >{\raggedright\arraybackslash}p{5cm}@{}}

\textbf{Tên viết tắt} & \textbf{Tên (English)} & \textbf{Diễn giải (Tiếng Việt)} \\
\midrule
\endfirsthead

BFS     & Breadth-First Search & Tìm kiếm theo chiều rộng \\
CMS     & Count-Min Sketch & Cấu trúc dữ liệu "sketch" để ước lượng tần suất phần tử trong luồng; cho ước lượng xấp xỉ với sai số có thể kiểm soát \\
CTDL    & Data Structures & Cấu trúc dữ liệu \\
DDoS    & Distributed Denial of Service & Tấn công từ chối dịch vụ phân tán \\
DSA     & Data Structures and Algorithms & Cấu trúc dữ liệu và giải thuật \\
eBPF    & extended Berkeley Packet Filter & Kỹ thuật cho phép chạy mã trong kernel để lọc/thu thập/xử lý gói mạng (eBPF) \\
IP      & Internet Protocol & Giao thức Internet dùng để định danh và chuyển gói tin \\
K8s     & Kubernetes & Hệ thống điều phối container (quản lý deployment, scaling, scheduling) \\
LPM     & Longest Prefix Match & Khớp tiền tố dài nhất \\
SCC     & Strongly Connected Component & Thành phần liên thông mạnh \\
SOC     & Security Operations Center & Trung tâm vận hành an ninh (giám sát và phản ứng sự cố) \\
YAML    & YAML Ain't Markup Language & Ngôn ngữ dữ liệu cấu hình dạng phân cấp (YAML) \\
MG      & Misra--Gries & Thuật toán Misra--Gries để xấp xỉ các phần tử có tần suất lớn bằng cách duy trì $k$ bộ đếm \\
CU      & Conservative Update & Biến thể Conservative Update của CMS \\
Mops    & Million operations per second & Triệu thao tác/giây \\
KB      & Kilobyte & 1 KB = 1024 bytes \\
CSV     & Comma-Separated Values & Tập tin dữ liệu phân tách dấu phẩy \\
L2      & Level-2 cache & Bộ nhớ đệm cấp 2 \\

MurmurHash3 & MurmurHash3 & Hàm băm không mật mã, nhanh và phổ biến để băm khóa trong hệ thống \\
Zipfian & Zipfian distribution & Phân phối Zipf (phân phối lệch; một vài phần tử rất phổ biến) \\
Heavy hitters & Heavy hitters & Các phần tử có tần suất cao trong luồng (phần tử 'nặng') \\
Count Sketch & Count Sketch & Biến thể sketch dùng để ước lượng tần suất, khác với CMS về cấu trúc và sai số \\
Single-pass & Single-pass algorithm & Thuật toán xử lý luồng trong một lần duyệt (mỗi phần tử chỉ xử lý một lần) \\
HashMap & HashMap (exact) & Bảng băm lưu trữ tần suất chính xác (đòi hỏi bộ nhớ tuyến tính) \\
Prometheus & Prometheus & Hệ thống giám sát và thu thập metric thời gian thực \\
Grafana & Grafana & Công cụ trực quan hóa và dashboard cho metrics \\
Pipeline & Pipeline & Dòng xử lý dữ liệu gồm các bước (collector → normalizer → aggregator → ...) \\
Collector & Collector & Thành phần thu thập sự kiện/thông điệp từ nguồn \\
Normalizer & Normalizer & Thành phần chuẩn hóa dữ liệu/sự kiện về dạng thống nhất \\
Aggregator & Aggregator & Thành phần tổng hợp dữ liệu từ nhiều nguồn/nút \\

FPR     & False Positive Rate & Tỷ lệ báo động giả (lưu lượng hợp lệ bị gắn cờ tấn công) \\
TCP     & Transmission Control Protocol & Giao thức truyền tải hướng kết nối, tin cậy \\
CNI     & Container Network Interface & Chuẩn cắm mạng cho container/Kubernetes (vd Calico, Cilium) \\
CI      & Confidence Interval & Khoảng tin cậy (thống kê), ở đây dùng mức 95\% \\
CO-RE   & Compile Once -- Run Everywhere & Kỹ thuật biên dịch eBPF một lần chạy trên nhiều phiên bản kernel nhờ BTF \\
BTF     & BPF Type Format & Định dạng thông tin kiểu dữ liệu của kernel, nền tảng cho CO-RE \\
kprobe  & Kernel probe & Điểm chèn (hook) động vào hàm kernel để chạy chương trình eBPF \\
NetworkPolicy & Kubernetes NetworkPolicy & Đối tượng K8s khai báo luật cho/chặn lưu lượng pod (do CNI thực thi) \\
MTTR    & Mean Time To Respond & Thời gian trung bình để phản hồi/xử lý sự cố \\
nmap    & Network Mapper & Công cụ quét cổng/dịch vụ mạng phổ biến \\
SLR     & Systematic Literature Review & Tổng quan tài liệu có hệ thống \\
SIEM    & Security Information and Event Management & Hệ quản lý thông tin và sự kiện an ninh (phân tích nhật ký) \\
NIDS    & Network Intrusion Detection System & Hệ phát hiện xâm nhập ở mức mạng \\
SRE     & Site Reliability Engineering & Kỹ nghệ độ tin cậy hệ thống (vận hành) \\
FR      & Functional Requirement & Yêu cầu chức năng \\
NFR     & Non-Functional Requirement & Yêu cầu phi chức năng \\
DLP     & Data Loss Prevention & Ngăn ngừa thất thoát dữ liệu (kiểm soát lối ra) \\
BFD     & Business Function Diagram & Sơ đồ phân rã chức năng nghiệp vụ \\
DFD     & Data Flow Diagram & Sơ đồ luồng dữ liệu \\
BPMN    & Business Process Model and Notation & Ký pháp mô hình hóa quy trình nghiệp vụ \\
ERD     & Entity--Relationship Diagram & Sơ đồ thực thể -- quan hệ (thiết kế CSDL) \\
UML     & Unified Modeling Language & Ngôn ngữ mô hình hóa thống nhất \\
UC      & Use Case & Ca sử dụng (trong phân tích yêu cầu) \\
C4      & C4 model & Mô hình kiến trúc phần mềm bốn mức (Context/Container/Component/Code) \\
RBAC    & Role-Based Access Control & Kiểm soát truy cập theo vai trò \\
OCI     & Open Container Initiative & Chuẩn ảnh và runtime container mở \\
GitOps  & GitOps & Mô hình triển khai khai báo, đồng bộ trạng thái cụm từ Git \\
RingBuf & Ring Buffer & Bộ đệm vòng của eBPF để chuyển sự kiện lên không gian người dùng \\
CSDL    & Database & Cơ sở dữ liệu \\

\end{longtable}

\clearpage

% ============================================================================
% Bang dien giai ky hieu
% ============================================================================

\chapter*{BẢNG DIỄN GIẢI KÝ HIỆU}
\addcontentsline{toc}{chapter}{BẢNG DIỄN GIẢI KÝ HIỆU}

\begin{longtable}{@{}p{3.5cm} p{10cm}@{}}

$S = (a_1, \ldots, a_m)$ & Luồng dữ liệu gồm $m$ phần tử \\
$U$                      & Không gian phần tử (universe), $|U| = n$ \\
$f_i$                    & Tần suất thực của phần tử $i$ trong luồng \\
$\hat{f}_x$              & Tần suất ước lượng của phần tử $x$ \\
$\varepsilon$            & Tham số sai số (error parameter) \\
$\delta$                 & Tham số xác suất thất bại (failure probability) \\
$F_k$                    & Frequency moment bậc $k$: $F_k = \sum f_i^k$ \\
$\|\mathbf{f}\|_1$       & Chuẩn $L_1$ của vector tần suất, $= F_1 = m$ \\
$w$                      & Chiều rộng (width) của ma trận CMS \\
$d$                      & Chiều sâu (depth) — số hàm băm \\
$h_i: U \to [w]$         & Hàm băm thứ $i$, ánh xạ phần tử sang cột \\
$G = (V, E)$             & Đồ thị có hướng với tập đỉnh $V$, tập cung $E$ \\

\end{longtable}

% Một số ký hiệu toán học bổ sung được dùng xuyên suốt báo cáo
\begin{longtable}{@{}p{3.5cm} p{10cm}@{}}

$m$                      & Độ dài luồng (số phần tử), nếu $S=(a_1,\ldots,a_m)$ \\
$k$                      & Tham số Misra--Gries (số bộ đếm $k$) \\
$F_0$                    & Số phần tử phân biệt trong luồng (distinct elements) \\
$M$                      & Ma trận Count--Min Sketch kích thước $d\times w$; $M[i][j]$ là bộ đếm \\
$O(\cdot)$               & Ký hiệu Big-O: cận trên về độ phức tạp thời gian/không gian \\
$\mathbb{E}[\cdot]$     & Kỳ vọng (expected value) \\
$\Pr[\cdot]$             & Xác suất (probability) \\
$e$                      & Hằng số Euler ($e\approx 2.71828$) \\
$\lceil \cdot \rceil$    & Hàm trần (ceiling) \\
$\ln$                    & Logarit tự nhiên \\
$\mathcal{H}$            & Họ hàm băm (family of hash functions) \\

\end{longtable}
